\bookchapter{第五十章：Kotlinはじめの一歩——Android開発の第一言語}

\chapterintro{この章で、いよいよ言語そのものを学びます。緊張はいりません。Kotlinは昔から「初心者にやさしい」ことで有名ですし、文法書を丸暗記する必要もありません。文法の2割を押さえるだけで、日常のコードの8割は読めて書けます。残りはプロジェクトを進めながら調べればよいのです。学習の原則を思い出してください。理解する＞暗記する、使える＞完璧、です。}

\booksection{変数と型：データに名前を付ける}

プログラムとは「データ＋ロジック」です。変数は、データに名前を付けて保存しておく仕組みのこと。Kotlinには2つの宣言方法があります。

\begin{lstlisting}[language=Kotlin]
val name = "たろう"        // val：読み取り専用の変数（デフォルトはこれがおすすめ）
var age = 18             // var：変更できる変数（あとから変えられる）

age = 19                 // <i class="fas fa-check-circle mr-1"></i> varは変えられる
// name = "はなこ"         // <i class="fas fa-times-circle mr-1"></i> valは変えられない。コンパイルエラー

val score: Int = 100     // 型を明示的に書いてもいい
val pi = 3.14            // 書かなくても自動で推論される
\end{lstlisting}

よく使う型は数えるほど。1枚の表で覚えてしまいましょう。

\begin{center}
\begin{tabularx}{\linewidth}{|X|X|X|}
\hline
\textbf{型} & \textbf{何を入れる？} & \textbf{例} \tabularnewline
\hline
Int & 整数 & 42 \tabularnewline
Double & 小数 & 3.14 \tabularnewline
String & 文字列 & "こんにちは" \tabularnewline
Boolean & 真/偽 & true / false \tabularnewline
List< T> & データの集まり & listOf(1, 2, 3) \tabularnewline
\hline
\end{tabularx}
\end{center}

\begin{center}\small\color{gray}Kotlinのよく使う基本型\end{center}

\begin{tcbinfo}{ワンポイント：val優先の原則}
Kotlinコミュニティの鉄則です。ふだんはvalを使い、本当に必要なときだけvarにする。valで足りる場所にvarを使うと、IDEが黄色い文字で警告してくれます。この習慣だけで「変数が意図せず書き換わる」系のバグのかなりの割合を防げます。コードの動きが少ないほど、頭の中で追いやすいのです。

\end{tcbinfo}

\booksection{関数：ロジックを「ボタン」にまとめる}

関数とは、名前のついたロジックのひとかたまりです。リモコンのボタンのようなもの——押す（呼び出す）と、一連の動作が実行されます。Kotlinではfunで宣言します。

\begin{lstlisting}[language=Kotlin]
fun greet(name: String): String {
    return "こんにちは、$name！"     // $記号で変数を文字列に差し込む
}

println(greet("たろう"))        // 出力：こんにちは、たろう！

// 1行の関数は式として簡潔に書ける
fun add(a: Int, b: Int) = a + b

println(add(3, 5))            // 出力：8
\end{lstlisting}

ポイントはいくつか。引数はかっこの中に書く（名前が先、型が後、コロンで区切る）。戻り値の型はかっこの後ろに書く。文字列テンプレートの\$変数名は、Kotlin名物の便利機能——Javaのように「文字列＋変数」とつなぎ合わせる必要はありません。

\booksection{条件とループ：プログラムに「考える」力をつける}

プログラムの賢さは2つの力から生まれます。判断できること（if/when）と、繰り返せること（for/while）です。

\begin{lstlisting}[language=Kotlin]
// 判断：if / else
val age = 17
if (age >= 18) {
    println("大人")
} else {
    println("未成年者")
}

// 多分岐：when（switchの百倍使いやすい）
val score = 85
val level = when {
    score >= 90 -> "優秀"
    score >= 60 -> "合格"
    else -> "その調子！"
}

// 繰り返し：forループ
val fruits = listOf("りんご", "バナナ", "オレンジ")
for (fruit in fruits) {
    println("$fruitが大好き")
}
\end{lstlisting}

whenはKotlinでいちばんよく使う多分岐の書き方です。Android開発では毎日お世話になります（たとえば状態に応じて別々の画面を出す、など）。

\booksection{Null安全性：Kotlinの「命綱」}

Javaの時代でいちばん有名なバグが、NullPointerException（ヌルポインタ例外）です。何も入っていない「空」のものにプログラムがアクセスして、その場でクラッシュします。統計上も、Androidのクラッシュランキングの常連1位です。Kotlinはこれを言語の仕組みレベルで解決しました。デフォルトではnullを許さないのです。

\begin{lstlisting}[language=Kotlin]
var name: String = "たろう"   // 普通の型：nullは絶対に許されない
// name = null                // <i class="fas fa-times-circle mr-1"></i> コンパイラが直接止めてくれる

var nickname: String? = null // 疑問符を付けると：nullを許す
nickname = "たーくん"

// null許容型を使うときは、安全に扱う必要がある：
println(nickname?.length)    // safe call：nullならnullを返すだけで、落ちない
val len = nickname?.length ?: 0   // Elvis演算子：nullのときはデフォルト値
println(len)
\end{lstlisting}

\begin{tcbwarn}{注意：?と?:を取り違えない}
String?は「nullを許す」、a ?: bは「aがnullならbを使う」という意味です（形がエルヴィスの髪型に似ているので、Elvis演算子と呼ばれます）。書き始めの頃にコンパイラから「null関連のエラー」が出たら、たいていはnull許容型の扱いを忘れているせいです。それは、言語があなたを守ってくれている合図なのです。

\end{tcbwarn}

\booksection{クラスとオブジェクト：「部品」を作る鋳型}

クラスはオブジェクトの「鋳型」で、オブジェクトはクラスから作った「製品」です。Androidの世界にあるもの——画面、ボタン、データ——は、すべてオブジェクトです。Kotlinにはdata class（データクラス）という便利な道具があり、データを「入れる」ことに特化していて、Javaなら数十行になるのを1行で済ませます。

\begin{lstlisting}[language=Kotlin]
// 普通のクラス：鋳型
class Dog(val name: String) {
    fun bark() = println("$name ワンワン！")
}

val dog = Dog("ポチ")      // オブジェクトを作る
dog.bark()                 // ポチ ワンワン！

// data class：データを入れる箱（equals/copyなど付き）
data class User(val id: Int, val name: String, val email: String)

val me = User(1, "たろう", "zhangsan@example.com")
val twin = me.copy(name = "たろう丸")   // コピーして1つのフィールドだけ変える
\end{lstlisting}

この先、おこづかい帳を作るときは、1件の記録がそのままdata classになります。data class Record(val amount: Double, val category: String, val date: String)。ほら、もう本物のプロジェクトのコード構造が書けるのです。

\booksection{ラムダと高階関数：Kotlinの「シンタックスシュガー」}

Androidのコードでは、\{ ... \}という「名前のない関数」を毎日見かけます。これがラムダ（匿名関数）です。ButtonのonClick = \{ count++ \}も1つのラムダで、「クリックされたら何をするか」をひと塊にして、ボタンに手渡しています。

\begin{lstlisting}[language=Kotlin]
// 普通の書き方：名前付き関数を定義してから渡す
fun double(x: Int) = x * 2
val result = listOf(1, 2, 3).map(::double)

// lambdaの書き方：その場で定義。手軽でストレート
val result2 = listOf(1, 2, 3).map { it * 2 }   // it = 現在の要素

// 絞り込み + 変換の連携（実際のプロジェクトで非常によくある形）
val bills = listOf(12.5, 88.0, 5.0, 200.0)
val bigBills = bills.filter { it > 50 }.sorted()
// [88.0, 200.0]——「50を超える請求を並べ替え」
\end{lstlisting}

使用頻度の高い「コレクションの連係技」を覚えましょう。filter（しぼり込み）、map（変換）、sorted（並べ替え）、sum（合計）。この先のリスト処理の90\%は、この4つの組み合わせです。Kotlinの哲学はこうです。1言で言い切れるロジックを、5言で書く必要はない。

\booksection{超初心者向けの文法早見カード}

1日で全部の文法を覚えようとしないこと。このカードを保存しておき、コードを書くときに照らし合わせれば、2週間もすれば自然と身につきます。

\begin{center}
\begin{tabularx}{\linewidth}{|X|X|}
\hline
\textbf{やりたいこと} & \textbf{書き方} \tabularnewline
\hline
読み取り専用の変数を定義 & \texttt{val x = 1} \tabularnewline
変更できる変数を定義 & \texttt{var x = 1} \tabularnewline
関数を定義 & fun 名前(引数): 戻り値の型 \{ ... \} \tabularnewline
条件分岐 & if (条件) \{...\} else \{...\} \tabularnewline
多分岐 & when \{ 条件 -> 結果 \} \tabularnewline
繰り返し & for (x in リスト) \{ ... \} \tabularnewline
文字列の組み立て & "こんにちは、\$name" \tabularnewline
nullの扱い & \texttt{x?.\allowbreak length ?: 0} \tabularnewline
データクラスを定義 & \texttt{data class X(val a: Int, val b: String)} \tabularnewline
\hline
\end{tabularx}
\end{center}

\begin{center}\small\color{gray}Kotlin文法早見カード（初心者版）\end{center}

\begin{tcbpractice}{やってみよう：ソフトを1つも入れずに、まずKotlinを10分遊ぶ}
ブラウザでplay.kotlinlang.org（Kotlin公式のオンライン練習場）を開き、この章のコードを順番に貼り付けて実行してみましょう。数字や文字を書き換えて、出力がどう変わるか観察します。この「遊んでから学ぶ」10分のおかげで、第五十一章に入ったときの負担が半分になります。

\end{tcbpractice}

\begin{tcbsummary}{この章のまとめ}
\begin{itemize}
  \item valを優先、varは慎重に。型は自動で推論されるが、コロンの後で明示してもよい
  \item 関数＝まとめたロジック。文字列テンプレートの\$はKotlin名物の便利機能
  \item 判断はif/when、繰り返しはfor。whenはAndroid開発の「何でも屋」
  \item Null安全性はKotlinの命綱。ふつうの型はnullを拒み、null許容型は?／?:での扱いが必須
  \item data classはデータを入れる箱。1フィールド1行で、本物のプロジェクトで毎日使う
  \item 文法の2割で仕事の8割をこなす。まず使えるようになり、それから極める
\end{itemize}

\end{tcbsummary}
